\subsection{Overview and Scope}
The research domain addressed by this thesis intersects five major fields: sensor-based condition monitoring and anomaly detection in industrial and IoT environments, deep learning architectures for time-series anomaly detection and predictive maintenance, quantum computing and quantum optimization, agentic AI and multi-agent autonomous systems, and combinatorial optimization for scheduling problems. This chapter provides a comprehensive and critical synthesis of the fifty collected papers, organized thematically to establish the scholarly context within which the QASAMAP framework makes its contributions. Each paper is engaged with directly, with its specific claims, methodologies, and findings cited and connected to the empirical evidence from the 100,000-reading smart manufacturing dataset analyzed in Chapter 5.

\textbf{Post-Audit Disclosure.} The literature review in this chapter situates QASAMAP within existing research by mapping the intended framework architecture to prior work. The actual Ver13 implementation differs from some of these mappings: the agentic layer is a sequential nine-agent LLM pipeline (Monitoring, Diagnosis, Planning, Reporting, Correlation, Explainability, Self-Reflection, Scheduling, Pattern Learning) coordinated by the \texttt{MultiAgentOrchestrator}, not a blackboard-based multi-agent system with Perception and Decision Agents. The ``quantum'' components are classical heuristics (\texttt{energy\_distance\_anomaly\_score}, \texttt{random\_constrained\_search}, \texttt{graph\_message\_passing\_score}) executed on the CPU, not variational quantum circuits on Qiskit Aer or D-Wave annealers. The reader should interpret forward-looking architectural claims in this chapter as design motivation rather than implemented functionality.

\subsection{Deep Learning for Anomaly Detection: Foundation and Advances} \cite{Biamonte2017}
The theoretical foundation for deep learning-based anomaly detection was established through a series of foundational surveys and comparative studies. Chalapathy and Chawla (2019) produced the most comprehensive survey systematically reviewing deep architectures including autoencoders, variational autoencoders, recurrent neural networks, convolutional networks, and generative adversarial networks across application domains spanning cyber-intrusion detection, fraud, medical imaging, and industrial monitoring. Their taxonomy distinguishes between unsupervised, semi-supervised, and supervised anomaly detection paradigms and demonstrates that reconstruction-based autoencoders trained exclusively on normal data are among the most effective approaches for industrial scenarios characterized by abundant normal data and rare anomalous events — precisely the condition exhibited in the 100,000-reading dataset where the 8.92\% anomaly rate provides overwhelmingly more normal than anomalous training examples \cite{Chalapathy2019}.

Building directly upon this autoencoder paradigm for IoT time-series environments, Gao et al. (2022) introduced TSMAE — the Memory-Augmented Autoencoder for IoT time series anomaly detection \cite{Gao2022}. TSMAE addresses a fundamental weakness of standard autoencoders: their tendency to generalize so effectively during training that they can reconstruct anomalous inputs with low error, undermining the anomaly detection principle. By augmenting the autoencoder's latent space with a fixed-size memory module that stores prototypical normal patterns and forcing the decoder to reconstruct inputs solely from memory-retrieved representations rather than directly from encoded latent vectors, TSMAE ensures that anomalous inputs — which do not match any stored normal prototype — generate high reconstruction errors. This memory augmentation mechanism directly addresses the false-negative problem in anomaly detection on the research dataset, where 8.92\% of readings are anomalous and where the five distinct failure modes create multiple potential normality-deviation patterns that a standard autoencoder might learn to reconstruct.

The specific challenge of multi-sensor anomaly detection in industrial IoT contexts was addressed by Zheng et al. (2023). This work investigates the detection of anomalies that manifest through the joint behavior of multiple sensor channels simultaneously — a scenario directly instantiated by the research dataset's five sensor modalities (temperature, vibration, humidity, pressure, energy consumption) where the correlation analysis reveals that temperature is the dominant anomaly predictor (r = 0.447) while other sensors contribute complementary discriminative information \cite{zhang_unsupervised_2021}. They demonstrate that deep architectures capable of modeling inter-sensor correlations substantially outperform single-sensor approaches, providing theoretical justification for the multi-modal encoder design in QASAMAP's CNN-BiLSTM backbone.

The correlation-based approach to multi-sensor anomaly detection was specifically investigated by Han et al. (2022) who demonstrated that anomalous machine states can be reliably detected by monitoring deviations from the learned normal cross-sensor correlation structure. This finding is particularly relevant to the research dataset, where the correlation matrix reveals that temperature and predicted RUL exhibit the strongest relationship (r = -0.193), and that disruptions to these inter-sensor relationships provide an early warning signal for impending failures \cite{Han2022}. The QASAMAP framework integrates correlation-based features as part of the statistical feature engineering pipeline described in Chapter 4.

\subsection{Convolutional and Recurrent Architectures for Time-Series Anomaly Detection}
The application of convolutional neural networks to time-series anomaly detection was advanced by Wen and Keyes (2019) who proposed a CNN-based segmentation approach that outperforms RNN-based methods on several benchmark datasets while achieving substantially lower computational cost \cite{wen_time_2019}. The transfer learning component of their approach — pre-training on source domain time series and fine-tuning on target domain data — is particularly relevant to the QASAMAP methodology, where the CNN branch of the multi-task backbone can be pre-trained on publicly available industrial sensor datasets (such as the NASA C-MAPSS turbofan dataset) before fine-tuning on the research dataset, improving generalization on minority failure classes with limited training examples.

The autocorrelation-based LSTM autoencoder enhances the standard LSTM autoencoder with explicit modeling of temporal autocorrelation patterns within sensor windows. By incorporating autocorrelation coefficients as additional input features, the model captures periodicities and repetitive patterns in normal machine operation that are disrupted during anomalous states \cite{homayouni_autocorrelation-based_2020}. This autocorrelation augmentation strategy directly informs the statistical feature engineering pipeline in QASAMAP, where autocorrelation coefficients are computed for each sensor channel within each sliding window and appended to the standard statistical features.

Network-level deep convolutional anomaly detection with residual learning was investigated by Chouhan, Khan, and Khan (2019). Their channel boosting mechanism, which applies multiple parallel convolutional branches with different kernel sizes to capture temporal patterns at multiple scales, directly parallels the QASAMAP CNN branch design that employs convolutional layers with kernel sizes of 3, 5, and 7. The residual connections proposed by Chouhan et al. prevent gradient vanishing in deep networks and are incorporated into the QASAMAP backbone as shortcut connections between adjacent convolutional layers \cite{Chouhan2019}.

GAN-based approaches to multivariate anomaly detection, particularly relevant for scenarios with limited labeled anomaly examples, were investigated by authors in "GAN-Based Anomaly Detection for Multivariate Time Series Using Polluted Training Set" published in IEEE Transactions on Knowledge and Data Engineering (VOL. 35, NO. 12). The polluted training set scenario — in which the training data contains an unknown fraction of anomalous examples — directly mirrors the research dataset, \textit{where the 8.92\% anomaly rate in the training data means that purely unsupervised models must be robust to a moderate degree of training set contamination. The GAN-based discriminator provides an additional anomaly score orthogonal to reconstruction error, and the ensemble of reconstruction error and discriminator scores is shown to achieve higher AUC-ROC than either component alone \cite{du_gan-based_2023}.}

\subsection{Real-Time and Industrial IoT Anomaly Detection}
The specific requirements of real-time anomaly detection in industrial IoT settings — including low-latency inference, resource-constrained edge deployment, and handling of streaming data. This framework, which combines lightweight LSTM-based anomaly detection with adaptive threshold adjustment based on operational context, is particularly relevant to QASAMAP's Monitoring Agent design: the requirement for one-minute cycle time inference aligns with the one-minute sensor sampling interval of the research dataset, and the adaptive threshold mechanism directly informs QASAMAP's machine-specific anomaly score calibration \cite{nizam_real-time_2022}.

The challenge of anomaly detection in smart city wireless sensor networks, analogous to the industrial sensor network studied here, was investigated by Jain, et al (2016). While this work addresses environmental rather than industrial sensors, its treatment of spatial correlation between co-located sensors and its robust handling of sensor failures versus genuine environmental anomalies provide methodological insights applicable to the multi-machine sensor network represented by the 50-machine fleet in the research dataset \cite{jain_anomaly_2016}.

The autoencoder-based framework for specifically distinguishing between process anomalies (genuine equipment degradation) and sensor failures (measurement errors). This distinction is critically important for the research dataset, where the zero missing values across all 1.3 million data points suggest high sensor reliability, but where real deployment environments will inevitably encounter sensor drift and failure events that must be distinguished from genuine machine degradation. QASAMAP's Monitoring Agent could incorporate a sensor health monitoring module inspired by this work \cite{lee_autoencoder-based_2024}.

Unsupervised anomaly detection for industrial robots using sliding-window convolutional variational autoencoders was investigated by Chen, et al, 2019), which demonstrates that the sliding window approach to converting continuous time series into fixed-length segments — exactly the approach adopted in QASAMAP's preprocessing layer — achieves state-of-the-art detection performance on robotic actuator degradation data. The variational autoencoder extension provides probabilistic anomaly scores with calibrated uncertainty estimates, directly motivating QASAMAP's uncertainty quantification approach in the anomaly detection head \cite{chen_unsupervised_2020}.

The comparative evaluation of unsupervised anomaly detection algorithms across multiple multivariate datasets, conducted by Goldstein and Uchida (2016) established Isolation Forest and Local Outlier Factor as the strongest classical baselines while identifying autoencoder-based methods as achieving the highest performance on high-dimensional, temporally structured data — directly supporting the QASAMAP design choice to adopt autoencoder-based detection over classical baselines \cite{Goldstein2016}.

\subsection{Quantum Computing: Foundations, Optimization, and Industrial Applications}
The quantum computing literature in the paper collection spans from introductory overviews to specialized algorithm development, providing a comprehensive foundation for the QASAMAP quantum components. They provide a systematic introduction to quantum software development, quantum programming languages, and the engineering challenges of deploying quantum algorithms on current hardware \cite{Epping2025}. The editorial highlights the critical importance of error mitigation, circuit optimization, and hybrid quantum-classical architectures — all of which are directly relevant to QASAMAP's quantum layer design.

The paradigm-shift characterization of quantum computing is articulated by Eldh and Sigrid, 2025, which argues that quantum computing represents not merely an incremental improvement over classical computation but a fundamentally new computational paradigm requiring rethinking of algorithm design from first principles. This perspective directly informs the QASAMAP quantum layer's approach of formulating the maintenance scheduling problem as a native quantum optimization problem (QUBO/Ising Hamiltonian) rather than applying quantum subroutines to a classically-designed algorithm \cite{eldh_quantum_2025}.

The practical application of quantum optimization to scheduling and routing problems — problems structurally analogous to the maintenance scheduling challenge in QASAMAP — is extensively documented by Rosendo, Fortunato, and Abreu (2025). They demonstrate that QAOA and quantum annealing achieve competitive solution quality compared to classical heuristics on vehicle routing problems of practical size on current NISQ hardware \cite{rosendo_quantum_2025}. Their conclusion that "quantum computing provides a promising approach to improving transportation optimization" extends directly to the maintenance scheduling domain, where the combinatorial structure of the scheduling problem mirrors that of vehicle routing.

The advanced quantum annealing approach to vehicle routing problems with time windows, developed by Holliday, Blount, Osaba, and Luu (2025), introduces constraint encoding techniques for handling time window constraints in quantum annealing formulations that are directly applicable to maintenance scheduling with shift-based time window constraints. Their QUBO formulation methodology — encoding both objective function and constraints as penalty terms in the Ising Hamiltonian — is adopted verbatim in QASAMAP's QAOA scheduling formulation \cite{holliday_advanced_2025}.

Tambunan et al. (2022) provide a detailed treatment of the QUBO formulation for routing problems with weighted edge segments, demonstrating that the D-Wave quantum annealer achieves solution quality within 5\% of optimal for problem sizes of up to 20 nodes — a result that supports the QASAMAP hierarchical decomposition strategy of partitioning the 50-machine fleet into 10-machine subgroups for QAOA optimization \cite{Tambunan2022}.

The comprehensive treatment of quantum optimization challenges and opportunities by Blekos et al. (2024) provides the theoretical framework within which QASAMAP's QAOA component is positioned. This paper systematically characterizes the complexity classes for which quantum optimization advantages are theoretically supported, identifies the barren plateau problem as the primary training challenge for variational quantum algorithms, and reviews mitigation strategies including layerwise training and warm-starting from classical solutions that are implemented in a potential future VQC training procedure \cite{Blekos2024}.

Practical quantum optimization applications bridging theory and hardware were explored by introducing software tools for mapping high-level optimization problems to quantum circuit implementations. The Quark framework's problem decomposition capabilities — partitioning large optimization problems into hardware-compatible subproblems — directly inspired QASAMAP's hierarchical fleet decomposition approach \cite{windgatter_quantum_2025}.

The digitized counterdiabatic quantum algorithms for logistics scheduling, developed by Dala, et al, 2024, introduce an enhanced QAOA variant that incorporates counterdiabatic driving terms to accelerate convergence and improve solution quality at low circuit depths. This enhancement is particularly valuable for NISQ hardware where circuit depth is limited by coherence time, and its application to logistics scheduling — a problem class that encompasses maintenance scheduling — directly supports the adoption of counterdiabatic QAOA as an enhancement of the standard QAOA formulation in QASAMAP \cite{dalal_digitized_2024}.

The practitioner's guide to quantum algorithms for optimization problems provides the algorithmic reference for QASAMAP's quantum optimization component, systematically comparing QAOA, quantum annealing, variational quantum eigensolver, and hybrid classical-quantum optimization approaches across problem classes defined by their combinatorial structure. The guide's recommendation of QAOA for scheduling problems with quadratic objective functions and linear constraints directly validates the QASAMAP scheduling formulation \cite{symons_practitioners_2023}.

The commercial applications of quantum computing, reviewed by Bova et al. (2021), provide industry-level evidence for the near-term applicability of quantum optimization to manufacturing and logistics problems. Their review identifies maintenance scheduling optimization as one of the three highest-value near-term application areas for quantum computing in industry — providing external validation for the commercial relevance of QASAMAP's quantum scheduling component \cite{Bova2021}.

The quantum machine learning library sQUlearn provides the Python implementation framework for the VQC components of QASAMAP, offering hardware-efficient ansatz designs, automatic gradient computation via the parameter-shift rule, and seamless integration with PennyLane's backend-agnostic quantum circuit execution engine. The sQUlearn documentation and benchmark results directly inform potential future VQC hyperparameter choices \cite{kreplin_squlearn_2025}.

The prediction of temporal variability in quantum processor error rates by Rodríguez-Soriano et al. (2025) addresses the practical challenge of scheduling quantum circuit execution on cloud quantum processors whose gate fidelity varies temporally due to environmental noise. This work motivates QASAMAP's adaptive quantum circuit scheduling strategy, which defers quantum computation to low-noise time windows predicted by the error rate forecasting model \cite{rodriguez-soriano_predicting_2025}.

\subsection{Hybrid Quantum-Classical and Quantum Machine Learning}
The hybrid quantum-classical reinforcement learning approach for industrial optimization in latent observation spaces, developed by Nagy, et al, 2025, demonstrates that encoding classical observation representations in quantum latent spaces enables reinforcement learning agents to discover policies that classical agents cannot find within practical training budgets. This finding motivates the integration of quantum feature encoding in QASAMAP's agentic pipeline, where maintenance policy optimization can potentially benefit from quantum-enhanced representation of the fleet's health state space \cite{nagy_hybrid_2025}.

The biological sequence comparison algorithm using quantum computers (Scientific Reports, 2023) provides a methodological case study of mapping combinatorial sequence alignment problems — structurally similar to the multi-machine maintenance schedule alignment problem — to quantum gate circuits. While the application domain differs from industrial maintenance, the QUBO formulation techniques and complexity analysis are directly transferable to the QASAMAP scheduling context.


\subsection{Agentic AI Systems: Architecture and Applications}
The paper collection contains the most current and comprehensive literature on agentic AI systems, providing a strong foundation for QASAMAP's agentic layer design. Sapkota, Roumeliotis, and Karkee (2026) in "AI Agents vs. Agentic AI: A Conceptual Taxonomy, Applications and Challenges" (Information Fusion, volume 126, article 103599) establish a rigorous conceptual distinction between individual AI agents (discrete, task-specific autonomous entities) and agentic AI systems (integrated architectures in which multiple agents collaborate, plan, reason, and execute over extended horizons). This distinction is fundamental to QASAMAP's multi-agent architecture: the individual agents are specialized LLM-powered modules, while the complete QASAMAP system — comprising nine agents coordinated by the MultiAgentOrchestrator in a sequential pipeline — constitutes an agentic AI system in the full sense defined by Sapkota et al \cite{sapkota_ai_2026}.

The comprehensive survey of agentic AI architectures by Abou Ali, Dornaika, and Charafeddine (2026) in Artificial Intelligence Review (DOI: 10.1007/s10462-025-11422-4) provides the most up-to-date synthesis of agentic AI design principles, identifying five key architectural components — perception modules, world models, reasoning engines, action executors, and memory systems — that characterize state-of-the-art agentic systems. QASAMAP's agentic layer implements all five components: the Monitoring and Diagnosis Agents serve as perception modules and world model updaters, the LLM reasoning provides the reasoning engine, the QAOA scheduler and work order generator constitute the action executor, and the ContextMemoryStore implements the memory system. This alignment confirms that QASAMAP's agentic design is consistent with the state-of-the-art architectural principles \cite{AbouAli2026}.

The rise of autonomous intelligent agents in the era of large language models, documented in "Agentic AI: The Rise of Autonomous Intelligent Agents in the Era of LLMs" (2025), establishes the LLM-powered reasoning core as the defining architectural innovation of contemporary agentic systems. While QASAMAP's initial implementation uses a structured reasoning framework rather than a full LLM backbone for computational efficiency reasons, the design is explicitly forward-compatible with LLM integration, and the Chapter 9 future directions explicitly identify LLM enhancement as a high-priority extension \cite{erukude_agentic_2025}.

The framework-level treatment of agentic AI architectures, protocols, and design challenges in "Agentic AI Frameworks: Architectures, Protocols, and Design Challenges" (arXiv:2508.10146, 2025) provides the protocol-level specifications for inter-agent communication, task delegation, and conflict resolution that inform QASAMAP's sequential agentic pipeline for agent coordination. The publish-subscribe communication pattern recommended by this work for multi-agent systems with heterogeneous agent types directly maps to QASAMAP's health assessment agentic pipeline \cite{derouiche_agentic_2025}.

The characterization of what agentic AI systems are and are not, presented in "Agentic AI Systems: What It Is and Isn't" (Global Business and Organizational Excellence, 2025), clarifies common misconceptions about agentic AI and establishes the key distinguishing properties: autonomous goal-directed behavior, persistent memory across interaction cycles, dynamic tool invocation, and self-reflection capabilities. QASAMAP's agentic pipeline implements all four properties: it pursues the persistent goal of fleet health optimization (autonomous goal-direction), maintains a continuously updated fleet health state in the ContextMemoryStore (persistent memory), invokes deep learning inference, quantum scheduling, and work order generation tools dynamically (dynamic tool invocation), and implements a performance monitoring module that evaluates its own decision outcomes (self-reflection)\cite{dwivedi_agentic_2026}.

The agentic retrieval-augmented generation survey (arXiv:2501.09136) provides the technical foundation for QASAMAP's integration of retrieval-based knowledge access in the agentic pipeline's reasoning process, enabling the agent to consult historical maintenance records, equipment manuals, and failure mode databases when formulating maintenance recommendations. The survey demonstrates that agentic RAG substantially improves the relevance and accuracy of agent-generated recommendations compared to parametric knowledge alone \cite{singh_agentic_2025}.


\subsection{Sensor Data Intelligence with LLMs and Deep QUestion Answering}
The application of large language models to sensor data understanding represents an emerging research frontier directly relevant to QASAMAP's interpretability objectives. LLMSense (arXiv:2403.19857, 2024) — "Harnessing LLMs for High-Level Reasoning over Spatiotemporal Sensor Traces" — demonstrates that LLMs can perform sophisticated causal reasoning about sensor data patterns when provided with structured natural language representations of sensor traces. This capability — interpreting why a temperature spike at a particular machine is anomalous given the concurrent readings from neighboring sensors and the machine's recent maintenance history — is precisely the type of contextual reasoning that QASAMAP's agentic pipeline must perform when generating human-interpretable maintenance recommendations \cite{ouyang_llmsense_2024}.

DeepSQA (DOI: 10.1145/3450268.3453529, 2021) — "Understanding Sensor Data via Question Answering" — introduces a deep learning framework for answering natural language questions about sensor data streams, enabling non-technical users to query machine health status in conversational terms. The DeepSQA architecture, which combines sensor feature encoding with BERT-based natural language understanding, provides a template for QASAMAP's human-agent interaction interface, where maintenance engineers can ask natural language questions about machine health and receive data-grounded answers \cite{singh_agentic_2025}.


\subsection{Combinatorial Optimization with Deep and Reinforcement Learning}
The application of deep reinforcement learning to dynamic vehicle routing with demand and traffic uncertainty (Operations Research Perspectives, 2025) provides methodological foundations for the adaptive optimization components of QASAMAP. The attention-based policy network proposed in this work — which learns to sequentially assign routes based on a learned representation of the current routing state — is architecturally analogous to the sequential maintenance assignment policy that QASAMAP's Planning Agent must learn, where machines must be assigned to maintenance slots based on a learned representation of the fleet health state.

The general modeling and simulation framework for dynamic problems (arXiv:2411.12406, 2024) provides a formal framework for representing the QASAMAP maintenance scheduling problem as a dynamic decision-making process with stochastic transitions, enabling the rigorous analysis of scheduling policy optimality that supports the comparison between QAOA and classical heuristic approaches in Chapter 8 \cite{horvath_general_2024}.